\subsection{LOCALLY COMPACT GROUPS OPERATING PROPERLY}

PROPOSITION 7. Let G be a locally compact group operating continuously on a Hausdorff space X. Then G operates properly on X if and only if, for each pair of points x, y of X, there is a neighbourhood \(V_{x}\) of x and a neighbourhood \(V_{y}\) of y such that the set K of all \(s \in G\) for which \(s \cdot V_{x} \cap V_{y} \neq \emptyset\) is relatively compact in G.

Let F be the compact space obtained by adjoining a point at infinity \(\omega\) to G, and let \(\Gamma\) be the graph of \(\rho: (s, x) \to s.x\) considered as a subset of \(F \times X \times X\) . Let us show that if the restriction of \(pr_{23}\) to \(\Gamma\) is proper, then \(\Gamma\) is closed in \(F \times X \times X\) . Indeed, this hypothesis implies that the map \(u: (t, s, x, y) \to (t, x, y)\) of \(F \times \Gamma\) into \(F \times X \times X\) is closed. If \(\Gamma'\) is the set of points \((s, s)\) in \(F \times G\) , where \(s \in G\) , then \(\Gamma'\) is closed in \(F \times G\) , since it is the graph of the canonical injection \(G \to F\) (Chapter I, §8, no. I, Corollary 2 to Proposition 2); hence the intersection \((\Gamma' \times X \times X) \cap (F \times \Gamma)\) is closed in \(F \times \Gamma\) , and it is immediately seen that its image under u is the set \(\Gamma\) considered as a subset of \(F \times X \times X\) ; hence \(\Gamma\) is closed in \(F \times X \times X\) . Now, we have \((\{\omega\} \times X \times X) \cap \Gamma = \emptyset\) . From the definition of F, therefore, for every point \((x, y) \in X \times X\) there is a neighbourhood W of \((x, y)\) in \(X \times X\) and a compact subset K of G such that

\[
((\mathbf {G} - \mathbf {K}) \times \mathbf {W}) \cap \Gamma ;
\]

is empty; and since we may take W to be a neighbourhood \(V_{x} \times V_{y}\) , where \(V_{x}\) and \(V_{y}\) are neighbourhoods of x and y respectively in X, the statement ``((G - K) × W) ∩ Γ = ∅'' becomes ``if s∉K, then s.Vx∩Vy=∅''. We have thus proved the necessity of the condition stated in the proposition. Conversely, suppose this condition is satisfied; let A be a set filtered by an ultrafilter \(\mathfrak{F}\), and let \(\alpha \to (s_{\alpha}, x_{\alpha})\) be a mapping of A into \(G \times X\) such that \(\lim_{\mathfrak{F}} x_{\alpha} = x\) and \(\lim_{\mathfrak{F}} s_{\alpha}\cdot x_{\alpha} = y\). Suppose that K, Vx and Vy satisfy the condition of the proposition. By hypothesis there is a set \(M \in \mathfrak{F}\) such that if \(\alpha \in M\) then \(x_{\alpha} \in V_{x}\) and \(s_{\alpha}\cdot x_{\alpha} \in V_{y}\), hence \(s_{\alpha} \in K\). This shows that \(\alpha \to s_{\alpha}\) converges with respect to \(\mathfrak{F}\), and the proof is complete.

If G is compact, the condition of Proposition 7 is trivially satisfied; we thus retrieve Proposition 2a).

Proposition 7 shows in particular that a discrete group G, operating continuously on a Hausdorff space X, operates properly on X if and only if, for each pair \((x, y)\) of points of X, there is a neighbourhood \(V_{x}\) of x and a neighbourhood \(V_{y}\) of y such that the set of points \(s \in G\) for which \(s \cdot V_{x} \cap V_{y} \neq \emptyset\) is finite.

PROPOSITION 8. Let G be a discrete group operating properly on a Hausdorff space X. Let x be a point of X and let \(K_{x}\) be the stabilizer of x. Then:

\begin{enumerate}
\def\labelenumi{\alph{enumi})}
\item
  The subgroup \(K_{x}\) is finite and there is an open set \(U \subset X\) , containing x, which is stable under \(K_{x}\) , and on which the equivalence relation induced by the relation defined by G is the equivalence relation defined by \(K_{x}\) .
\item
  The canonical mapping \(\mathbf{U} / \mathbf{K}_x\to \mathbf{X} / \mathbf{G}\) is a homeomorphism of \(\mathbf{U} / \mathbf{K}_x\) onto an open neighbourhood of the class of \(x\) in \(\mathbf{X} / \mathbf{G}\) .
\end{enumerate}

By Proposition 7, \(\mathbf{K}_x\) is finite. To construct an open set \(\mathbf{U}\) which satisfies the required conditions, notice first that by Proposition 7 there is an open set \(\mathrm{U}_0\) containing \(x\) and such that the set \(\mathbf{K}\) of \(s \in \mathbf{G}\) for which \(s. \mathrm{U}_0 \cap \mathrm{U}_0 \neq \emptyset\) is finite. Clearly \(\mathbf{K}_x \subset \mathbf{K}\) ; let \(s_1, \ldots, s_n\) be the elements of \(\mathbf{K} - \mathbf{K}_x\) . If we put \(x_i = s_i.x\) ( \(1 \leqslant i \leqslant n\) ), then \(x_i \neq x\) for each \(i\) ; since \(\mathbf{X}\) is Hausdorff, there is for each index \(i\) an open neighbourhood \(\mathrm{V}_i\) of \(x\) and an open neighbourhood \(\mathrm{V}_i'\) of \(s_i.x\) such that \(\mathrm{V}_i \cap \mathrm{V}_i' = \emptyset\) . Let \(\mathrm{U}_i = \mathrm{V}_i \cap s_i^{-1}. \mathrm{V}_i'\) ; then \(\mathrm{U}_i\) is clearly open and contains \(x\) , and we have \(\mathrm{U}_i \cap s_i. \mathrm{U}_i \subset \mathrm{V}_i \cap \mathrm{V}_i' = \emptyset\) . Let \(\mathrm{U}' = \mathrm{U}_0 \cap \mathrm{U}_1 \cap \ldots \cap \mathrm{U}_n\) ; \(\mathrm{U}'\) is open, contains \(x\) and is such that \(\mathrm{U}' \cap s. \mathrm{U}' = \emptyset\) for \(s \notin \mathbf{K}_x\) . Putting \(\mathrm{U} = \bigcap_{t \in \mathbf{K}_x} t. \mathrm{U}'\) we finally obtain an open set, stable under \(\mathbf{K}_x\) , containing \(x\) and such that \(\mathrm{U} \cap s. \mathrm{U} = \emptyset\) for \(s \notin \mathbf{K}_x\) : \(\mathrm{U}\) is the open set required.

The fact that the canonical mapping \(\mathbf{U} / \mathbf{K}_x\to \mathbf{X} / \mathbf{G}\) is a homeomorphism of \(\mathbf{U} / \mathbf{K}_x\) onto an open set in \(\mathbf{X} / \mathbf{G}\) follows from Chapter I, § 5, no. 2, Proposition 4, since \(\mathbf{U}\) is open and the equivalence relation defined by \(\mathbf{G}\) is open (§ 2, no. 4, Lemma 2).

COROLLARY. If we suppose in addition that \(K_{x}=\{e\}\) , then the point x has an open neighbourhood U such that the restriction to U of the canonical mapping \(X\to X/G\) is a homeomorphism of U onto an open subset of X/G.

